\makeindex

\newcommand{\minitab}[2][l]{\begin{tabular}{@{}#1@{}}#2\end{tabular}}


\definecolor{ListingBG}{rgb}{0.91,0.91,0.91}
\definecolor{shadecolor}{rgb}{0.92,0.92,0.92}

\hyphenation{Open-SHMEM}

\renewcommand{\chaptername}{Chapter}
\renewcommand{\appendixname}{Annex}
\Crefname{appendix}{Annex}{Annexes}

% Place some penalty for doing the break
\def\gb{\penalty10000\hskip 0pt plus 8em\penalty4800\hskip 0pt plus-8em%
\penalty10000}

% Searchable underscore with allowed line break after it
\DeclareRobustCommand{\_}{\texttt{\char`\_}\discretionary{}{}{}}

\def\colorswapnt{\colorlet{saved}{.}\color{ForestGreen}}
\def\colorswapot{\colorlet{saved}{.}\color{red}}
\def\prevcolor{\color{saved}}

\newcommand{\newtext}[1]{\textcolor{ForestGreen}{#1}}
\newcommand{\oldtext}[1]{\textcolor{magenta}{\sout{#1}}}
\newcommand{\insertDocVersion}{0.1}
\newcommand{\openshmem}[1][]{%
  {Open\-SHMEM\ifthenelse{\equal{#1}{}}{}{~#1}}\xspace}
\newcommand{\HEADER}[1]{\textit{#1}}
\newcommand{\FUNC}[1]{\textit{#1}}
\newcommand{\CTYPE}[1]{\textit{#1}}
\newcommand{\VAR}[1]{\textit{#1}}
\newcommand{\ENVVAR}[1]{\textit{#1}}
\newcommand{\CONST}[1]{\textit{#1}}
\newcommand{\KEYWORD}[1]{\textit{#1}}
%%
\newcommand{\CorCpp}{\textit{C/C++}\xspace}
\newcommand{\Cstd}[1][]{%
  \textit{C\ifthenelse{\equal{#1}{}}{}{#1}}\xspace}
%%
\newcommand{\TYPE}{\emph{TYPE}}
\newcommand{\TYPENAME}{\emph{TYPENAME}}
\newcommand{\SIZE}{\emph{SIZE}}

\newcommand{\source}{\textit{source}}
\newcommand{\dest}{\textit{dest}}
\newcommand{\PUT}{\textit{Put}}
\newcommand{\GET}{\textit{Get}}
\newcommand{\OPR}[1]{\textit{#1}}
\newcommand{\shmemprefix}{\textit{SHMEM\_}}
\newcommand{\shmemprefixLC}{\textit{shmem\_}}
\newcommand{\shmemgprefix}{\textit{shmemg\_}}
\newcommand{\ith}{${\textit{i}^{\text{\tiny th}}}$}
\newcommand{\jth}{${\textit{j}^{\text{\tiny th}}}$}

%% Generate indexed reference.
\newcommand{\EnvVarIndex}[1]{\index{#1}}
\newcommand{\FuncIndex}[1]{\index{#1}}
\newcommand{\LibConstIndex}[1]{\index{#1}}
\newcommand{\LibHandleIndex}[1]{\index{#1}}
\newcommand{\TableIndex}[1]{\index{#1}\index{Tables!#1}}
%% Write text and generate reference.
\newcommand{\EnvVarRef}[1]{\ENVVAR{#1}\EnvVarIndex{#1}}
\newcommand{\FuncRef}[1]{\FUNC{#1}\FuncIndex{#1}}
\newcommand{\LibConstRef}[1]{\CONST{#1}\LibConstIndex{#1}}
\newcommand{\LibHandleRef}[1]{\CONST{#1}\LibHandleIndex{#1}}
\newcommand{\TableCaptionRef}[1]{\caption{#1}\TableIndex{#1}}
%% Specialized declaration/creation and generate reference.
\newcommand{\EnvVarDecl}[1]{\EnvVarRef{#1}}
\newcommand{\FuncDecl}[1]{{\ListingsCurrentStyle{#1}}\FuncIndex{#1}}
\newcommand{\FuncParam}[1]{{\ListingsParamStyle{#1}}}
\newcommand{\LibConstDecl}[2][\CorCpp]{%
  \parbox[t]{5cm}{~\\[-4pt] #1: \\\hspace*{8mm} \LibConstRef{#2} \\~}}
\newcommand{\LibHandleDecl}[2][\CorCpp]{%
  \parbox[t]{0pt}{~\\[-4pt] #1: \\\hspace*{8mm} \LibHandleRef{#2} \\~}}

\newcommand{\subsubsubsection}[1]{\paragraph{#1}\mbox{}\par\nobreak}
\Crefname{paragraph}{Section}{Sections}
\setlength{\parindent}{0pt}

\begin{acronym}
\acro{RMA}{Remote Memory Access}
\acro{RMO}{\emph{Remote Memory Operation}}
\acro{AMO}{Atomic Memory Operation}
\acro{PE}{\emph{Processing Element}}
\acrodefplural{PE}[PEs]{\emph{Processing Elements}}
\acro{PGAS}{\emph{Partitioned Global Address Space}}
\acro{API}{\emph{Application Programming Interface}}
\acro{GPU}{\emph{Graphics Processing Unit}}
\end{acronym}

% Remove whitespace and retain commas when using cleveref's `\cref`.
\makeatletter
\let\org@@cref\@cref
\renewcommand*{\@cref}[2]{%
  \edef\process@me{%
    \noexpand\org@@cref{#1}{\zap@space#2 \@empty}%
  }\process@me
}
\makeatother
\newcommand{\ChangelogRef}[1]{\\See \cref{#1}.\space}

% Grab current listings style for use in environment escape to LaTeX.
\makeatletter
\newcommand\ListingsCurrentStyle{}
\lst@AddToHook{Output}{\global\let\ListingsCurrentStyle\lst@thestyle}
\lst@AddToHook{OutputOther}{\global\let\ListingsCurrentStyle\lst@thestyle}
\makeatother

%
% Use Sans Serif font for sections
%
\makeatletter
\def\section{\@startsection {section}{1}{\z@}{-3.5ex plus -1ex minus
-.2ex}{2.3ex plus .2ex}{\Large\sf}}
\def\subsection{\@startsection{subsection}{2}{\z@}{-3.25ex plus -1ex minus
-.2ex}{1.5ex plus .2ex}{\large\sf}}
\def\subsubsection{\@startsection{subsubsection}{3}{\z@}{-3.25ex plus
-1ex minus -.2ex}{1.5ex plus .2ex}{\normalsize\sf\bf}}
\def\paragraph{\@startsection {paragraph}{4}{\z@}{3.25ex plus 1ex minus .2ex}
{-1em}{\normalsize\sf\bf}}
\makeatother

%
% Listings styles
%

\newcommand{\ListingsParamStyle}{\bfseries\itshape\color{CadetBlue}}

\definecolor{DarkGreen}{rgb}{0, 0.25, 0}

\lstdefinestyle{SourceListingsDefault}{
  language={C},
  breaklines=true,
  breakatwhitespace=true,
  escapechar=@,
  tabsize=2,
  basicstyle=\ttfamily\footnotesize,
  showstringspaces=false,
  frame=single,
  classoffset=0,
  keywords={\#include, \#pragma, \#if, \#ifdef, \#else, \#endif},
  keywordstyle=\bfseries\color{CadetBlue},
  classoffset=1,
  keywords={
    if, else, for, while, do, return, break, continue,
    static, const, volatile,
    sizeof, typedef, struct, union,
    switch, case, default, goto,
  },
  keywordstyle=\bfseries\color{Orchid},
  classoffset=2,
  keywords={
    __device__, _Thread_local, _Noreturn,
    void, bool, _Bool, unsigned, char, short, int, long,
    float, double, _Complex, complex, size_t, ptrdiff_t,
    int8_t, int16_t, int32_t, int64_t,
    uint8_t, uint16_t, uint32_t, uint64_t,
    shmem_ctx_t, shmem_team_t, shmem_team_config_t,
    FILE,
  },
  keywordstyle=\color{DarkGreen},
  classoffset=3,
  keywords={NULL},
  keywordstyle=\color{MidnightBlue},
  classoffset=4,
  keywords={TYPE, TYPENAME, SIZE},
  keywordstyle=\ListingsParamStyle,
  commentstyle=\itshape\color{BrickRed},
  stringstyle=\color{DarkOrchid},
  backgroundcolor=\color{white},
}

\definecolor{gray}{rgb}{0.97,0.97,0.97}

\lstdefinestyle{PrototypeListingsDefault}{
  style=SourceListingsDefault,
  frame=none,
  backgroundcolor=\color{gray},
}

\lstset{style=SourceListingsDefault}

%
% Library API description template commands
%

\newcommand{\apisummary}[1]{
    #1
\hfill
}

\newenvironment{apidefinition}{
\begin{description}
\item[SYNOPSIS] \hfill \\ \\
\vspace{-2em}
}
{
\end{description} % ends \begin{description} from \apidescription
\end{description}
}

\lstnewenvironment{Csynopsis}
{
  \textbf{C:}
  \lstset{style=PrototypeListingsDefault}
}{}

\lstnewenvironment{CsynopsisCol}
{
  \lstset{style=PrototypeListingsDefault}
}{}

\newenvironment{apiarguments}{
\newcommand{\apiargument}[3]{
\begin{tabular}{p{2cm} p{2cm} p{10cm}}
\textbf{##1} & \textit{##2} & {##3} \\
\end{tabular}
}
\hfill
\item[DESCRIPTION] \hfill

\begin{description}
\item[Arguments] \hfill \\
}
{
\hfill
\end{description}
}

\newcommand{\apidescription}[1]{
\begin{description}
\vspace{-1em}
\item[API Description] \hfill \\
    \sloppy
    #1
\hfill
}

\newcommand{\apireturnvalues}[1]{
\item[Return Values] \hfill \\
    #1
}

\newcommand{\apinotes}[1]{
\item[Notes] \hfill \\
    #1
}
